When we evaluated an open-weight LLM on model-in-the-loop adversarial NLP benchmarks (ANLI), we found a counterintuitive result: calibration \textit{improved} relative to the clean baseline ($\Delta$ECE $= -0.041$ for ANLI R1). The adversarial benchmark, designed to expose model failures, made the model appear better calibrated. Yet human-adversarially crafted NLI examples (AdvGLUE MNLI) produced the opposite --- a 7.1 percentage-point ECE increase ($\Delta$ECE $= +0.071$). Calibration robustness under adversarial conditions is not universal: it depends critically on \textit{how} adversarial examples were constructed. We measure Expected Calibration Error (ECE) directly on adversarial NLP benchmark splits for an open-weight LLM (Llama-2-7b-hf), finding that calibration degradation is conditional on adversarial construction method and task type. RLHF alignment moderates this pattern in a benchmark-type-dependent manner --- consistently improving calibration robustness on model-in-loop adversarial conditions (100\% moderation rate across all ANLI rounds) while exacerbating degradation on static human-adversarial benchmarks (alignment tax: $\Delta\Delta$ECE $= -0.026$ on AdvGLUE). In this single-model pilot study, these findings motivate adversarial calibration measurement as a distinct research problem from adversarial accuracy robustness, and provide a conditional framework --- construction method $\times$ task type $\times$ RLHF alignment --- for understanding when and how adversarial perturbation may degrade model reliability signals in deployment.
